\bookchapter{부록 A: 새 스마트폰 개봉 “파워풀” 절차}

\booksection{개봉 및 기기 점검 “3부작”:}

\begin{itemize}
  \item \textbf{외관 확인}: 기기 본체, 화면에 스크래치, 찌그러짐이 있는지 점검.
  \item \textbf{기능 테스트}: 전화를 걸어 이어피스와 마이크를 테스트, 사진을 찍어 카메라 점검.
\end{itemize}

\booksection{활성화 및 초기 설정 “핵심 리마인더”:}

\begin{itemize}
  \item \textbf{네트워크}: 먼저 안정적인 Wi-Fi에 연결.
  \item \textbf{계정}: Apple ID, 화웨이 계정, 샤오미 계정 등 로그인.
  \item \textbf{기초 설정}: "내 스마트폰 찾기" 켜기, 신뢰할 수 있는 잠금 화면 비밀번호 설정.
  \item \textbf{eSIM 설정}: 스마트폰이 eSIM을 지원하는 경우(iPhone, 일부 Android 플래그십), 설정에서 eSIM 요금제를 추가할 수 있으며, 실물 SIM 카드 삽입 불필요. 해외 출국 시 eSIM이 특히 편리, 현지 통신사의 eSIM 요금제를 직접 구매 가능.
  \item \textbf{데이터 마이그레이션}: 구형 스마트폰에서 데이터 이전 시, 두 대의 스마트폰 배터리가 충분한지 확인. iPhone은 "빠른 시작", Android는 공식 교체 앱 사용.
\end{itemize}

\booksection{앱 설치 “보안 우선”:}

\begin{itemize}
  \item \textbf{공식 스토어 우선}: Apple은 App Store, Android는 스마트폰 자체 앱 스토어 사용.
  \item \textbf{권한 함부로 부여 금지}: 설치 시 팝업되는 권한 요청, "이 앱이 정말 필요한가?" 생각.
\end{itemize}
